% =====================================================================
%  Chapter 3 - Target Tracking Pipeline
% =====================================================================

\chapter{Target Tracking Pipeline}
\label{ch:tracking}

This chapter is focused on describing the main characteristics of the target
tracking module. The objective of the algorithm is to estimate, at every instant, the state
of the opponents around the ego vehicle, i.e.\ their position, speed, heading and
acceleration, starting from the detections produced by the perception stage of each
sensor. As anticipated in Section~\ref{sec:sw_arch}, the list of the opponents with their
estimated state is the output of the perception module and becomes the input of the
planning module.

The estimation of the state of each opponent follows the structure of an Extended Kalman
Filter (EKF), the latter chosen to handle the nonlinearities of the dynamic model that
describes the motion of the vehicle. The filter alternates two essential steps: the
\emph{prediction}, in which the state is propagated in time according to a Constant
Curvature and Acceleration (CCA) model, and the \emph{update}, in which the predicted
state is corrected when a new measurement becomes available. Around the filter, the
target tracking module manages all the operations required to work with real
detections: the association of each measurement with the corresponding opponent, the
handling of the measurements that arrive out of sequence (OOSM), and the management of
the tracks, i.e.\ their creation, confirmation and deletion.

\begin{figure}[H]
    \centering
    \resizebox{\textwidth}{!}{% Block diagram of the target tracking pipeline (Chapter 3).
\begin{tikzpicture}[
    font=\footnotesize,
    >={Stealth[length=2mm]},
    blk/.style={draw, rounded corners=2pt, align=center, minimum height=10mm, inner sep=3pt},
    io/.style={blk, minimum width=18mm, fill=gray!8},
    mod/.style={blk, minimum width=20mm, fill=blue!6},
    ekf/.style={blk, minimum width=20mm, fill=blue!18, very thick},
    grp/.style={draw, dashed, rounded corners=4pt, inner sep=5pt},
    lab/.style={font=\scriptsize, align=center, text=black!70},
    line/.style={->, thick}
]
% main chain
\node[mod] (oosm) {OOSM\\handling};
\node[io, left=20mm of oosm] (perc) {Perception\\pipelines};
\node[ekf, right=9mm of oosm] (pred) {Prediction\\(CCA model)};
\node[mod, right=9mm of pred] (assoc) {Data association\\(Mahalanobis,\\Munkres)};
\node[ekf, right=12mm of assoc] (upd) {Update\\(correction)};
\node[blk, minimum width=16mm, right=9mm of upd] (tracks) {Tracks};
\node[mod, below=9mm of upd] (mgmt) {Track management\\(creation, confirmation,\\deletion)};

% external inputs and output
\node[io, above=9mm of perc] (raw) {Sensors\\(LiDAR, RaDAR,\\camera)};
\node[io, above=9mm of pred] (db) {Track\\database};
\node[io, above=9mm of assoc] (ego) {Ego state\\(estimation)};
\node[io, above=9mm of tracks] (plan) {Planning};

% data flow
\draw[line] (raw) -- node[lab, right]{raw data} (perc);
\draw[line] (perc) -- node[lab, above, pos=0.4]{detections} (oosm);
\draw[line] (oosm) -- node[lab, above]{$\mathbf{z},\,t_s$} (pred);
\draw[line] (pred) -- (assoc);
\draw[line] (assoc) -- node[lab, above]{assigned} (upd);
\draw[line] (assoc.south) |- node[lab, above, pos=0.75]{unassigned} (mgmt.west);
\draw[line] (upd) -- (tracks);
\draw[line] (mgmt.east) -| (tracks.south);
\draw[line] (db) -- node[lab, right]{curvature $\kappa$} (pred);
\draw[line] (ego) -- node[lab, right]{ego pose at $t_s$} (assoc);
\draw[line] (tracks) -- node[lab, right]{opponents} (plan);

% feedback: the tracks are the starting point of the next prediction
\coordinate (fb) at ($(mgmt.south)+(0,-5mm)$);
\coordinate (fbr) at ($(tracks.east)+(5mm,0)$);
\draw[line] (tracks.east) -- (fbr) |- (fb) -| node[lab, above, pos=0.25]{$\hat{\mathbf{x}},\,\mathbf{P}$ of the tracks} (pred.south);

\node[grp, fit=(oosm)(pred)(assoc)(upd)(tracks)(mgmt)(fb)(fbr), label={[lab]below:Target tracking}] {};
\end{tikzpicture}
}
    \caption{Main steps of the target tracking pipeline.}
    \label{fig:tt_pipeline}
\end{figure}

Following the same pipeline; Section~\ref{sec:tt_prediction}
describes the prediction step of the EKF. Section~\ref{sec:tt_association} explains how
each measurement is associated with the corresponding track, leveraging on the Mahalanobis
distance gating and the Munkres algorithm, and how the measurements that arrive out of sequence
are handled. Section~\ref{sec:tt_correction} presents the measurement models and the
update step, in which the predicted state is corrected with the associated measurement.
Section~\ref{sec:tt_management} is about tracks management. Finally,
Section~\ref{sec:tt_baseline} overviews the main performance limits of the current
pipeline, highlighting estimate degradation in rapidly changing dynamic conditions, i.e.\ the starting point of the research work.

% =====================================================================
\section{Prediction}
\label{sec:tt_prediction}

The prediction step propagates the state of each track, together with its uncertainty,
from the time of the last update to the time of the next measurement, according to a
model of the motion of the opponent.

\subsection{Prediction Step}
\label{subsec:tt_prediction}

The EKF describes the evolution of the state of the opponent and its relation with the
measurements through the discrete-time nonlinear system
\begin{equation}
    \begin{aligned}
        \mathbf{x}_{k+1} &= f_d(\mathbf{x}_k) + \mathbf{w}_k, \\
        \mathbf{z}_k &= h(\mathbf{x}_k) + \mathbf{v}_k,
    \end{aligned}
    \label{eq:tt_system}
\end{equation}
where $f_d(\cdot)$ is the motion model, which propagates the state $\mathbf{x}_k$ over the
interval between two filter iterations, and $h(\cdot)$ is the measurement model, which
depends on the sensor that produced the measurement $\mathbf{z}_k$
(Section~\ref{sec:tt_correction}). The process noise $\mathbf{w}_k$ represents the
confidence in the motion model, i.e.\ how reliable the prediction is expected to be,
while the measurement noise $\mathbf{v}_k$ represents the accuracy of the measurement; both are modelled as zero-mean Gaussian white noises, with covariance
matrices $\mathbf{Q}_k$ and $\mathbf{R}_k$ respectively.

Since $f_d$ is nonlinear, it is linearized around the last estimate,
\begin{equation}
    \mathbf{F}_k = \left.\frac{\partial f_d}{\partial \mathbf{x}}\right|_{\hat{\mathbf{x}}_k}.
    \label{eq:tt_F_def}
\end{equation}
The estimated state $\hat{\mathbf{x}}$ is associated with a measure of its uncertainty,
represented by the covariance matrix $\mathbf{P}$ of the filter. Given the estimate
$\hat{\mathbf{x}}_k$ and its covariance $\mathbf{P}_k$, the prediction step computes the a
priori estimate and covariance at the next iteration,
\begin{equation}
    \begin{aligned}
        \hat{\mathbf{x}}_{k+1|k} &= f_d(\hat{\mathbf{x}}_k), \\
        \mathbf{P}_{k+1|k} &= \mathbf{F}_k\,\mathbf{P}_k\,\mathbf{F}_k^\top + \mathbf{Q}_k,
    \end{aligned}
    \label{eq:tt_prediction}
\end{equation}
where the covariance is propagated forward with the motion model and increased by the
contribution of the process noise.

In the target tracking module the prediction step is performed either when a new
measurement is received, in which case the state is predicted to the validity time of
the measurement, or when a given time has elapsed since the last iteration, so that the
estimate is kept up to date also in the absence of measurements. Since the measurements
come from different sensors, with different update rates and latencies, the interval
$\Delta t = t_{k+1} - t_k$ between two consecutive predictions is not constant. For this
reason the discrete-time equations used to propagate the state and the process noise are
derived through an exact discretization of the continuous-time model, which is valid for
any value of $\Delta t$ and is described in Section~\ref{subsec:tt_cca}.

\subsection{Constant Curvature and Acceleration Model}
\label{subsec:tt_cca}

The motion of the opponent is described by a Constant Curvature and Acceleration (CCA)
model. The state of each opponent is defined as
\begin{equation}
    \renewcommand{\arraystretch}{1.3}\setlength{\jot}{8pt}
    \mathbf{x} = \begin{bmatrix} x & y & v & a & \psi \end{bmatrix}^\top,
    \label{eq:tt_state}
\end{equation}
where $(x,\,y)$ is the position of the opponent, $v$ its speed, $a$ its longitudinal
acceleration and $\psi$ its heading.

\paragraph{Continuous-time model.} The CCA model assumes that, over a prediction
interval, the opponent travels along a path of constant curvature $\kappa$ with a constant
longitudinal acceleration. The continuous-time dynamics are
\begin{equation}
    \renewcommand{\arraystretch}{1.3}\setlength{\jot}{8pt}
    \left\{
    \begin{aligned}
        \dot x(t)    &= v(t)\cos\psi(t), \\
        \dot y(t)    &= v(t)\sin\psi(t), \\
        \dot v(t)    &= a(t), \\
        \dot a(t)    &= 0, \\
        \dot\psi(t)  &= \kappa\,v(t),
    \end{aligned}
    \right.
    \label{eq:tt_cca}
\end{equation}
where the yaw rate is expressed as $\dot\psi = \kappa\,v$: once the path is fixed, the
rotation of the vehicle depends only on how fast it travels along it. In compact form,
the system is written as $\dot{\mathbf{x}}(t) = f(\mathbf{x}(t),\kappa)$.

The a priori knowledge of the trajectories that the opponent may follow, i.e.\ the racing
lines stored in the track database introduced in Section~\ref{sec:sw_arch}, is exploited
as a prior and modelled as an external input: the curvature $\kappa$ is not estimated by
the filter, but read from the database. At every prediction the point closest to the estimated position of the opponent
is searched on the optimal trajectory, starting from the index found at the previous step;
if its distance from the opponent exceeds \SI{0.5}{m}, the search is extended to all the
other trajectories of the database, and the curvature of the closest point is used for
the prediction.

\paragraph{Discrete-time model.} Starting from the continuous-time
model~\eqref{eq:tt_cca}, the discrete-time motion model is obtained through its exact
discretization,
\begin{equation}
    \renewcommand{\arraystretch}{1.3}\setlength{\jot}{8pt}
    \mathbf{x}_{k+1} = \mathbf{x}_k
    + \int_{t_k}^{t_{k+1}} f(\mathbf{x}(\tau),\kappa)\,\mathrm{d}\tau.
    \label{eq:tt_exact_disc}
\end{equation}
Under the assumption of constant curvature and acceleration over the interval, the
integrals can be solved in closed form. For acceleration, speed and heading,
\begin{equation}
    \renewcommand{\arraystretch}{1.3}\setlength{\jot}{8pt}
    \begin{aligned}
        a_{k+1}    &= a_k + \int_{t_k}^{t_{k+1}} 0\,\mathrm{d}\tau = a_k, \\
        v_{k+1}    &= v_k + \int_{t_k}^{t_{k+1}} a_k\,\mathrm{d}\tau = v_k + a_k\,\Delta t, \\
        \psi_{k+1} &= \psi_k + \int_{t_k}^{t_{k+1}} \kappa\,v(\tau)\,\mathrm{d}\tau
                    = \psi_k + \kappa\left(v_k\,\Delta t + \tfrac{1}{2}\,a_k\,\Delta t^2\right),
    \end{aligned}
    \label{eq:tt_cca_vapsi}
\end{equation}
while, since $\mathrm{d}\psi = \kappa\,v\,\mathrm{d}\tau$, the integrals of the position
reduce to integrals in the heading,
\begin{equation}
    \renewcommand{\arraystretch}{1.3}\setlength{\jot}{8pt}
    \begin{aligned}
        x_{k+1} &= x_k + \int_{t_k}^{t_{k+1}} v(\tau)\cos\psi(\tau)\,\mathrm{d}\tau
                 = x_k + \frac{\sin\psi_{k+1} - \sin\psi_k}{\kappa}, \\
        y_{k+1} &= y_k + \int_{t_k}^{t_{k+1}} v(\tau)\sin\psi(\tau)\,\mathrm{d}\tau
                 = y_k - \frac{\cos\psi_{k+1} - \cos\psi_k}{\kappa}.
    \end{aligned}
    \label{eq:tt_cca_xy}
\end{equation}
The discrete-time motion model is therefore
\abovedisplayskip=16pt plus 3pt minus 3pt \belowdisplayskip=16pt plus 3pt minus 3pt
\abovedisplayshortskip=16pt plus 3pt minus 3pt \belowdisplayshortskip=16pt plus 3pt minus 3pt
\begin{equation}
    \renewcommand{\arraystretch}{1.3}\setlength{\jot}{8pt}
    f_d(\mathbf{x}_k) =
    \begin{bmatrix}
        x_k + \dfrac{\sin\psi_{k+1} - \sin\psi_k}{\kappa} \\[2ex]
        y_k - \dfrac{\cos\psi_{k+1} - \cos\psi_k}{\kappa} \\[2ex]
        v_k + a_k\,\Delta t \\
        a_k \\
        \psi_k + \kappa\left(v_k\,\Delta t + \tfrac{1}{2}\,a_k\,\Delta t^2\right)
    \end{bmatrix}.
    \label{eq:tt_cca_discrete}
\end{equation}
When the curvature tends to zero, the opponent travels along a straight line and the
model reduces to
\abovedisplayskip=16pt plus 3pt minus 3pt \belowdisplayskip=16pt plus 3pt minus 3pt
\abovedisplayshortskip=16pt plus 3pt minus 3pt \belowdisplayshortskip=16pt plus 3pt minus 3pt
\begin{equation}
    \renewcommand{\arraystretch}{1.3}\setlength{\jot}{8pt}
    \lim_{\kappa \to 0} f_d(\mathbf{x}_k) =
    \begin{bmatrix}
        x_k + \left(v_k\,\Delta t + \tfrac{1}{2}\,a_k\,\Delta t^2\right)\cos\psi_k \\
        y_k + \left(v_k\,\Delta t + \tfrac{1}{2}\,a_k\,\Delta t^2\right)\sin\psi_k \\
        v_k + a_k\,\Delta t \\
        a_k \\
        \psi_k
    \end{bmatrix},
    \label{eq:tt_cca_straight}
\end{equation}
which is used when $|\kappa| < \SI{e-4}{m^{-1}}$ to avoid the division by zero
in~\eqref{eq:tt_cca_discrete}.

Applying~\eqref{eq:tt_F_def} to~\eqref{eq:tt_cca_discrete}, the linearized state matrix
used in the prediction of the covariance is
\abovedisplayskip=16pt plus 3pt minus 3pt \belowdisplayskip=16pt plus 3pt minus 3pt
\abovedisplayshortskip=16pt plus 3pt minus 3pt \belowdisplayshortskip=16pt plus 3pt minus 3pt
\begin{equation}
    \renewcommand{\arraystretch}{1.3}\setlength{\jot}{8pt}
    \mathbf{F}_k =
    \begin{bmatrix}
        1 & 0 & \Delta t\cos\psi_{k+1} & \tfrac{1}{2}\Delta t^2\cos\psi_{k+1} & \dfrac{\cos\psi_{k+1} - \cos\psi_k}{\kappa} \\[2ex]
        0 & 1 & \Delta t\sin\psi_{k+1} & \tfrac{1}{2}\Delta t^2\sin\psi_{k+1} & \dfrac{\sin\psi_{k+1} - \sin\psi_k}{\kappa} \\[2ex]
        0 & 0 & 1 & \Delta t & 0 \\
        0 & 0 & 0 & 1 & 0 \\
        0 & 0 & \kappa\,\Delta t & \tfrac{1}{2}\kappa\,\Delta t^2 & 1
    \end{bmatrix},
    \label{eq:tt_F}
\end{equation}
while, in the straight-line case~\eqref{eq:tt_cca_straight}, it becomes
\abovedisplayskip=16pt plus 3pt minus 3pt \belowdisplayskip=16pt plus 3pt minus 3pt
\abovedisplayshortskip=16pt plus 3pt minus 3pt \belowdisplayshortskip=16pt plus 3pt minus 3pt
\begin{equation}
    \renewcommand{\arraystretch}{1.3}\setlength{\jot}{8pt}
    \mathbf{F}_k =
    \begin{bmatrix}
        1 & 0 & \Delta t\cos\psi_k & \tfrac{1}{2}\Delta t^2\cos\psi_k & -\left(v_k\,\Delta t + \tfrac{1}{2}\,a_k\,\Delta t^2\right)\sin\psi_k \\
        0 & 1 & \Delta t\sin\psi_k & \tfrac{1}{2}\Delta t^2\sin\psi_k & \phantom{-}\left(v_k\,\Delta t + \tfrac{1}{2}\,a_k\,\Delta t^2\right)\cos\psi_k \\
        0 & 0 & 1 & \Delta t & 0 \\
        0 & 0 & 0 & 1 & 0 \\
        0 & 0 & 0 & 0 & 1
    \end{bmatrix}.
    \label{eq:tt_F_straight}
\end{equation}
All the quantities in~\eqref{eq:tt_F} and~\eqref{eq:tt_F_straight} are evaluated at the
last estimate $\hat{\mathbf{x}}_k$.

\paragraph{Process noise.} Two sources of uncertainty are considered. The first is the
variation of the acceleration, which the model assumes constant: a white noise $w_a$ with
spectral density $q_a$ acts on its derivative, so that the acceleration evolves as a
random walk and can follow the braking and traction phases. The second is the
uncertainty of the curvature read from the database: since the set of trajectories that
the opponent may follow is used as prior information, this uncertainty is lower, but it
accounts for the small differences between the actual trajectory of the opponent and the
one prescribed by the database. A white perturbation $w_\kappa$ of the curvature, with
spectral density $q_\kappa$, produces an error $v\,w_\kappa$ on the yaw rate. The
continuous-time model with process noise becomes
\abovedisplayskip=16pt plus 3pt minus 3pt \belowdisplayskip=16pt plus 3pt minus 3pt
\abovedisplayshortskip=16pt plus 3pt minus 3pt \belowdisplayshortskip=16pt plus 3pt minus 3pt
\begin{equation}
    \renewcommand{\arraystretch}{1.3}\setlength{\jot}{8pt}
    \dot{\mathbf{x}}(t) = f(\mathbf{x}(t),\kappa) + \mathbf{G}\,\mathbf{w}(t),
    \qquad
    \mathbf{G} =
    \begin{bmatrix}
        0 & 0 & 0 & 0 & v \\
        0 & 0 & 0 & 1 & 0
    \end{bmatrix}^\top,
    \qquad
    \mathbf{Q}_c = \begin{bmatrix} q_\kappa & 0 \\ 0 & q_a \end{bmatrix},
    \label{eq:tt_noise}
\end{equation}
where $\mathbf{w} = [w_\kappa,\,w_a]^\top$ and $\mathbf{Q}_c$ is its spectral density
matrix. The discrete-time process noise covariance is obtained through the exact
discretization of the noise, i.e.\ by integrating its effect over the prediction interval
through the linearized dynamics,
\abovedisplayskip=16pt plus 3pt minus 3pt \belowdisplayskip=16pt plus 3pt minus 3pt
\abovedisplayshortskip=16pt plus 3pt minus 3pt \belowdisplayshortskip=16pt plus 3pt minus 3pt
\begin{equation}
    \renewcommand{\arraystretch}{1.3}\setlength{\jot}{8pt}
    \mathbf{Q}_k = \int_0^{\Delta t} e^{\mathbf{A}\tau}\,\mathbf{G}\,\mathbf{Q}_c\,
        \mathbf{G}^\top e^{\mathbf{A}^\top\tau}\,\mathrm{d}\tau,
    \label{eq:tt_Q_def}
\end{equation}
where $\mathbf{A}$ is the Jacobian of the continuous-time dynamics~\eqref{eq:tt_cca},
evaluated at the last estimate,
\abovedisplayskip=16pt plus 3pt minus 3pt \belowdisplayskip=16pt plus 3pt minus 3pt
\abovedisplayshortskip=16pt plus 3pt minus 3pt \belowdisplayshortskip=16pt plus 3pt minus 3pt
\begin{equation}
    \renewcommand{\arraystretch}{1.3}\setlength{\jot}{8pt}
    \mathbf{A} = \left.\frac{\partial f}{\partial \mathbf{x}}\right|_{\hat{\mathbf{x}}_k} =
    \begin{bmatrix}
        0 & 0 & \cos\psi & 0 & -v\sin\psi \\
        0 & 0 & \sin\psi & 0 & \phantom{-}v\cos\psi \\
        0 & 0 & 0 & 1 & 0 \\
        0 & 0 & 0 & 0 & 0 \\
        0 & 0 & \kappa & 0 & 0
    \end{bmatrix}.
    \label{eq:tt_A}
\end{equation}
The integral is solved analytically, and the resulting closed-form expressions of the
elements of $\mathbf{Q}_k$, which depend on $v$, $\psi$, $\kappa$ and $\Delta t$, are
evaluated at every prediction. The result is a full $5 \times 5$ matrix: although the
noise acts directly only on the acceleration and on the heading, it propagates
indirectly to the other states through the integration of the dynamics. The two spectral
densities $q_\kappa$ and $q_a$ are the tuning parameters of the model, and their ratio
with respect to the measurement noise sets how much the filter relies on the prediction
rather than on the measurements.

% =====================================================================
\section{Data Association}
\label{sec:tt_association}

\subsection{Mahalanobis Distance}
\label{subsec:tt_mahalanobis}

\subsection{Munkres Algorithm}
\label{subsec:tt_munkres}

\subsection{Out-of-Sequence Measurements}
\label{subsec:tt_oosm}

% =====================================================================
\section{Correction}
\label{sec:tt_correction}

\subsection{Measurement Model}
\label{subsec:tt_meas_model}

\subsection{Measurement Noise}
\label{subsec:tt_meas_noise}

\subsection{Update}
\label{subsec:tt_update}

% =====================================================================
\section{Track Management}
\label{sec:tt_management}

% =====================================================================
\section{Baseline Performance}
\label{sec:tt_baseline}

\FloatBarrier
